\section{NVIDIA Jetson}

\subsection{Jetson Nano}

Jetson Nano đại diện cho nhóm thiết bị Edge AI có tài nguyên hạn chế. Thiết bị
phù hợp để đánh giá mức tối thiểu mà model có thể vận hành, đồng thời làm rõ ảnh
hưởng của bộ nhớ, power mode và thermal throttling. Phiên bản JetPack, CUDA,
TensorRT, dung lượng RAM và chế độ công suất phải được ghi theo đúng thiết bị
thực nghiệm, vì các biến thể phần cứng có thể khác nhau.

\subsection{Jetson Orin Nano}

Jetson Orin Nano đại diện cho thế hệ GPU nhúng mới hơn, có kiến trúc và năng lực
tăng tốc khác Jetson Nano. Thiết bị được dùng để đánh giá khả năng mở rộng hiệu
năng của cùng pipeline model, đồng thời kiểm tra tính di động của cơ chế triển
khai và Edge Runtime giữa hai lớp node.

\subsection{So sánh tài nguyên phần cứng}

\begin{table}[H]
\centering
\caption{Thông số hai thiết bị Jetson (theo đặc tả của NVIDIA \cite{jetson})}
\label{tab:jetson-hardware}
\begin{tabularx}{\textwidth}{L{3.6cm}YY}
\toprule
\textbf{Thành phần} & \textbf{Jetson Nano 4GB} & \textbf{Jetson Orin Nano 8GB}\\
\midrule
Kiến trúc GPU & NVIDIA Maxwell & NVIDIA Ampere\\
Nhân GPU & 128 CUDA core, không có Tensor Core & 1024 CUDA core, 32 Tensor Core thế hệ 3\\
CPU & 4 nhân ARM Cortex-A57, 1,43 GHz & 6 nhân ARM Cortex-A78AE\\
RAM & 4 GB LPDDR4 64-bit, 25,6 GB/s & 8 GB LPDDR5 128-bit, 68 GB/s\\
Năng lực AI & 472 GFLOPS (FP16) & 40 TOPS INT8 (67 TOPS ở chế độ Super)\\
JetPack mới nhất & 4.6.x (L4T R32.7) & 6.x (L4T R36)\\
CUDA / TensorRT & 10.2 / 8.2 & 12.x / 8.6--10.x\\
Power mode & 5 W, 10 W (MAXN) & 7 W, 15 W, 25 W (Super)\\
Hỗ trợ INT8 & Không có đơn vị INT8 chuyên dụng & Tensor Core hỗ trợ INT8\\
\bottomrule
\end{tabularx}
\end{table}

Khác biệt quan trọng nhất đối với tối ưu mô hình là Jetson Nano không có Tensor
Core và phần cứng không hỗ trợ tăng tốc INT8, nên trên thiết bị này lợi ích chủ
yếu đến từ FP16 và từ việc giảm FLOPs bằng cắt tỉa; còn trên Jetson Orin Nano,
INT8 qua Tensor Core là nguồn tăng tốc chính. Jetson Nano cũng dừng ở JetPack 4.6
(Ubuntu 18.04, Python 3.6 mặc định), trong khi Orin Nano dùng JetPack 6
(Ubuntu 22.04); điều này buộc image runtime phải build riêng cho từng loại thiết bị.
Phiên bản phần mềm và power mode cụ thể khi đo được ghi lại trong
Bảng~\ref{tab:experiment-devices}.

\subsection{CUDA và GPU Acceleration}

CUDA cung cấp mô hình lập trình và runtime để thực thi kernel trên GPU NVIDIA.
Trong pipeline, TensorRT sử dụng CUDA để lập lịch kernel và truyền dữ liệu. Tối
ưu end-to-end cần hạn chế sao chép CPU--GPU, tái sử dụng bộ đệm, dùng stream
phù hợp và đồng bộ chỉ tại những điểm cần thiết.

GPU acceleration không tự động bảo đảm latency thấp. Kích thước input, batch,
memory allocation, camera decode và hậu xử lý trên CPU có thể trở thành nút thắt.
Vì vậy, đồ án đo từng giai đoạn của pipeline và tổng thời gian đầu cuối.